\documentclass[a4paper,12pt]{article}
\usepackage[utf8]{inputenc}
\usepackage[T1]{fontenc}
\usepackage[ngerman]{babel}
\usepackage{graphicx}
\usepackage{setspace}
\usepackage{geometry}
\usepackage{titling}
\usepackage{lmodern}
\usepackage{amsmath}
\usepackage[style=ieee,sorting=none]{biblatex}
\addbibresource{Literaturverzeichnis.bib}
\usepackage{fancyhdr} % Paket für benutzerdefinierte Header
\usepackage{titlesec} % Paket zum Anpassen der Kapitelüberschrift
\usepackage{parskip}
\usepackage{float}
\usepackage{hyperref}
\usepackage{makecell}
\usepackage{listings}
\usepackage{xcolor}
\usepackage{multirow}
\usepackage{tabularx} 
\usepackage{wasysym}
\usepackage{longtable}
\usepackage{booktabs}
\usepackage{array}
\usepackage{calc}

\geometry{left=3cm,right=3cm,top=3cm,bottom=3cm}

% Anpassen der Kapitelüberschrift (größer und dicker)
\titleformat{\section}[block]
  {\normalfont\bfseries\fontsize{30pt}{36pt}\selectfont} % Schriftart und Größe (fett und sehr groß)
  {\thesection}{1em}{} % Nummerierung und Abstand
\titlespacing*{\section}{0pt}{-1cm}{1cm} % Abstand: 3cm vor der Überschrift und 1cm danach

% Anpassen der Unterkapitelüberschrift (subsection)
\titleformat{\subsection}[block]
  {\normalfont\bfseries\Large} % Schriftart und Größe (fett und größer)
  {\thesubsection}{1em}{} % Nummerierung und Abstand
\titlespacing*{\subsection}{0pt}{0.5cm}{0.5cm} % Abstand: 0.5cm vor der Überschrift und 0.5cm danach

\titleformat{\subsubsection}[block]
  {\normalfont\bfseries\large} % Schriftart und Größe (fett und größer)
  {\thesubsubsection}{1em}{} % Nummerierung und Abstand
\titlespacing*{\subsubsection}{0pt}{0.5cm}{0.5cm} % Abstand: 0.5cm vor der Überschrift und 0.5cm danach

% Fancyhdr für benutzerdefinierte Header
\pagestyle{fancy}
\fancyhf{} % Alle Felder zurücksetzen

% Kapitel und Abschnitt im Header
\fancyhead[L]{\nouppercase{\rightmark}} % Kapitelnummer + Kapitelname links
% Header für Kapitel und Abschnitte definieren
\renewcommand{\sectionmark}[1]{\markright{\thesection\quad #1}} % Abschnitt mit Nummer
\renewcommand{\subsectionmark}[1]{\markright{\thesubsection\quad #1}} % Unterabschnitt mit Nummer
\begin{document}
\onehalfspacing
\fancyfoot[C]{\thepage}
% Logo
\begin{center}
\thispagestyle{empty}
    \includegraphics[width=0.4\textwidth]{RWU Logo.png}\hfill
    
\includegraphics[width=0.4\textwidth]{arc42-logo.png} \\[1cm]
    \Large \textsc{Hochschule Ravensburg-Weingarten} \\[1cm]
    \rule{\textwidth}{0.4pt} \\[0.4cm]
    {\huge \textbf{Software-Architektur Arc42 Dokumentation}} \\[0.1cm]
    \rule{\textwidth}{0.4pt} \\[1cm]

    \noindent
    \begin{minipage}[t]{0.35\textwidth}
        \begin{flushleft} \large
        \textit{Verfasser:} \\
        Philipp Söllner 35601
        Tom Petzold 34917
        Sascha Vusatjuk 35650
        \end{flushleft}
    \end{minipage}
    \hfill
    \begin{minipage}[t]{0.5\textwidth}
        \begin{flushleft} \large
        \textit{Betreuer:} \\
        Prof. Dr. Sebastian Mauser \\
        \makebox[1.1\linewidth][l]{Hochschule Ravensburg-Weingarten} \\
        \end{flushleft}
    \end{minipage}
    \vfill
    Datum
\end{center}
\pagebreak
\pagenumbering{roman}
\fancyfoot[C]{\thepage}


\thispagestyle{plain} % Inhaltsverzeichnis ohne Header
\tableofcontents
\pagebreak

\pagenumbering{arabic}
\setcounter{page}{1}

\section{Einführung und Ziele}

\subsection{Aufgabenstellung}
{\large \textbf{Inhalt}}\\
Der Bewerbungstracker ist eine webbasierte Anwendung zur
Verwaltung und Nachverfolgung von Bewerbungen. Nutzer können Bewerbungen anlegen, deren Status verwalten und den Überblick über ihren Bewerbungsprozess behalten.

Die \textbf{Kernfunktionalitäten} des Bewerbungstrackers sind:

\begin{itemize}
\item Bewerbungen anlegen mit (Firma, Position, Standort usw.)
\item Status der Bewerbungen ändern
\item Bewerbungen auf einem Kanban-Board visualisieren und per Drag \& Drop verschieben
\item Statistiken basierend auf allen Bewerbungen und Verlaufseinträgen einsehen
\end{itemize}

{\large \textbf{Motivation}}\\
Das Projekt ist im Rahmen der Lehrveranstaltung Software-Architektur entstanden. Es dient als Demonstration einer sauberen Schichtenarchitektur und soll den Studierenden einen Einblick geben, wie diese richtig implementiert wird und worauf dabei zu achten ist.

\subsection{Qualitätsziele}

offen

\subsection{Stakeholder}\label{_stakeholder}

{\large \textbf{Inhalt}}
\begin{longtable}[]{@{}
  >{\raggedright\arraybackslash}p{(\columnwidth - 4\tabcolsep) * \real{0.22}}
  >{\raggedright\arraybackslash}p{(\columnwidth - 4\tabcolsep) * \real{0.30}}
  >{\raggedright\arraybackslash}p{(\columnwidth - 4\tabcolsep) * \real{0.48}}@{}}
\toprule\noalign{}
\textbf{Rolle} & \textbf{Person / Gruppe} & \textbf{Erwartungshaltung} \\
\midrule\noalign{}
\endhead
\bottomrule\noalign{}
\endlastfoot
Entwickler & Tom Petzold, Sascha Vusatjuk, Philipp Söllner & Klare, saubere Schichtentrennung und lauffähiges Front-/Backend \\\\
Prüfer & Sebastian Mauser & Nachvollziehbare und saubere Schichtenarchitektur, sowie eine vollständige arc42-Dokumentation \\\\ 
Endnutzer & Jobsuchende & Übersichtliche Benutzeroberfläche zum anlegen von Bewerbungen und einfache Bedienung\\
\end{longtable}
\newpage
\section{Randbedingungen}

\begin{longtable}[]{@{}
  >{\raggedright\arraybackslash}p{(\columnwidth - 4\tabcolsep) * \real{0.44}}
  >{\raggedright\arraybackslash}p{(\columnwidth - 4\tabcolsep) * \real{0.56}}@{}}
\toprule\noalign{}
\textbf{Randbedingungen} & \textbf{Erläuterung} \\
\midrule\noalign{}
\endhead
\bottomrule\noalign{}
\endlastfoot
Schichtenarchitektur muss verwendet werden &  Die Anwendung muss in der Schichtenarchitektur umgesetzt werden, da im Rahmen der Vorlesung ein Projekt mit einer selbstausgewählten Software-Architektur umgesetzt werden muss \\\\
Arc42 als Dokumentationsform & Das Projekt ist auf Vorgabe des Professors gemäß dem arc42-Template zu dokumentieren \\\\
Begrenzte Ressourcen und Zeit & Das Projekt wird von drei Studierenden neben dem Studium umgesetzt und muss bis zum 29.06.2026 fertiggestellt werden \\
\end{longtable}
\newpage
\section{Kontextabgrenzung}\label{section-context-and-scope}

\textbf{Inhalt}

Die Kontextabgrenzung grenzt das System gegen alle Kommunikationspartner
(Nachbarsysteme und Benutzerrollen) ab. Sie legt damit die externen
Schnittstellen fest und zeigt damit auch die Verantwortlichkeit (scope)
Ihres Systems: Welche Verantwortung trägt das System und welche
Verantwortung übernehmen die Nachbarsysteme?

Differenzieren Sie fachlichen (Ein- und Ausgaben) und technischen
Kontext (Kanäle, Protokolle, Hardware), falls nötig.

\textbf{Motivation}

Die fachlichen und technischen Schnittstellen zur Kommunikation gehören
zu den kritischsten Aspekten eines Systems. Stellen Sie sicher, dass Sie
diese komplett verstanden haben.

\textbf{Form}

Verschiedene Optionen:

\begin{itemize}
\item
  Diverse Kontextdiagramme
\item
  Listen von Kommunikationsbeziehungen mit deren Schnittstellen
\end{itemize}

\textbf{Weiterführende Informationen}

Siehe \href{https://docs.arc42.org/section-3/}{Kontextabgrenzung} in der
online-Dokumentation (auf Englisch!).

\subsection{Fachlicher Kontext}\label{_fachlicher_kontext}

\textbf{Inhalt}

Festlegung \textbf{aller} Kommunikationsbeziehungen (Nutzer, IT-Systeme,
\ldots\hspace{0pt}) mit Erklärung der fachlichen Ein- und Ausgabedaten
oder Schnittstellen. Zusätzlich (bei Bedarf) fachliche Datenformate oder
Protokolle der Kommunikation mit den Nachbarsystemen.

\textbf{Motivation}

Alle Beteiligten müssen verstehen, welche fachlichen Informationen mit
der Umwelt ausgetauscht werden.

\textbf{Form}

Alle Diagrammarten, die das System als Blackbox darstellen und die
fachlichen Schnittstellen zu den Nachbarsystemen beschreiben.

Alternativ oder ergänzend können Sie eine Tabelle verwenden. Der Titel
gibt den Namen Ihres Systems wieder; die drei Spalten sind:
Kommunikationsbeziehung, Eingabe, Ausgabe.

\textbf{\textless Diagramm und/oder Tabelle\textgreater{}}

\textbf{\textless optional: Erläuterung der externen fachlichen
Schnittstellen\textgreater{}}

\subsection{Technischer Kontext}\label{_technischer_kontext}

\textbf{Inhalt}

Technische Schnittstellen (Kanäle, Übertragungsmedien) zwischen dem
System und seiner Umwelt. Zusätzlich eine Erklärung (\emph{mapping}),
welche fachlichen Ein- und Ausgaben über welche technischen Kanäle
fließen.

\textbf{Motivation}

Viele Stakeholder treffen Architekturentscheidungen auf Basis der
technischen Schnittstellen des Systems zu seinem Kontext.

Insbesondere bei der Entwicklung von Infrastruktur oder Hardware sind
diese technischen Schnittstellen durchaus entscheidend.

\textbf{Form}

Beispielsweise UML Deployment-Diagramme mit den Kanälen zu
Nachbarsystemen, begleitet von einer Tabelle, die Kanäle auf
Ein-/Ausgaben abbildet.

\textbf{\textless Diagramm oder Tabelle\textgreater{}}

\textbf{\textless optional: Erläuterung der externen technischen
Schnittstellen\textgreater{}}

\textbf{\textless Mapping fachliche auf technische
Schnittstellen\textgreater{}}
\newpage
\section{Lösungsstrategie}\label{section-solution-strategy}

\textbf{Inhalt}

Kurzer Überblick über die grundlegenden Entscheidungen und
Lösungsansätze, die Entwurf und Implementierung des Systems prägen.
Hierzu gehören:

\begin{itemize}
\item
  Technologieentscheidungen
\item
  Entscheidungen über die Top-Level-Zerlegung des Systems,
  beispielsweise die Verwendung gesamthaft prägender Entwurfs- oder
  Architekturmuster,
\item
  Entscheidungen zur Erreichung der wichtigsten Qualitätsanforderungen
  sowie
\item
  relevante organisatorische Entscheidungen, beispielsweise für
  bestimmte Entwicklungsprozesse oder Delegation bestimmter Aufgaben an
  andere Stakeholder.
\end{itemize}

\textbf{Motivation}

Diese wichtigen Entscheidungen bilden wesentliche „Eckpfeiler`` der
Architektur. Von ihnen hängen viele weitere Entscheidungen oder
Implementierungsregeln ab.

\textbf{Form}

Fassen Sie die zentralen Entwurfsentscheidungen \textbf{kurz} zusammen.
Motivieren Sie, ausgehend von Aufgabenstellung, Qualitätszielen und
Randbedingungen, was Sie entschieden haben und warum Sie so entschieden
haben. Vermeiden Sie redundante Beschreibungen und verweisen Sie eher
auf weitere Ausführungen in Folgeabschnitten.

\textbf{Weiterführende Informationen}

Siehe \href{https://docs.arc42.org/section-4/}{Lösungsstrategie} in der
online-Dokumentation (auf Englisch!).
\newpage
\section{Bausteinsicht}\label{section-building-block-view}

\textbf{Inhalt}

Die Bausteinsicht zeigt die statische Zerlegung des Systems in Bausteine
(Module, Komponenten, Subsysteme, Klassen, Schnittstellen, Pakete,
Bibliotheken, Frameworks, Schichten, Partitionen, Tiers, Funktionen,
Makros, Operationen, Datenstrukturen, \ldots\hspace{0pt}) sowie deren
Abhängigkeiten (Beziehungen, Assoziationen, \ldots\hspace{0pt})

Diese Sicht sollte in jeder Architekturdokumentation vorhanden sein. In
der Analogie zum Hausbau bildet die Bausteinsicht den
\emph{Grundrissplan}.

\textbf{Motivation}

Behalten Sie den Überblick über den Quellcode, indem Sie die statische
Struktur des Systems durch Abstraktion verständlich machen.

Damit ermöglichen Sie Kommunikation auf abstrakterer Ebene, ohne zu
viele Implementierungsdetails offenlegen zu müssen.

\textbf{Form}

Die Bausteinsicht ist eine hierarchische Sammlung von Blackboxen und
Whiteboxen (siehe Abbildung unten) und deren Beschreibungen.

\begin{figure}
\centering
\includegraphics{images/05_building_blocks-DE.png}
\caption{}
\end{figure}

\textbf{Ebene 1} ist die Whitebox-Beschreibung des Gesamtsystems,
zusammen mit Blackbox-Beschreibungen der darin enthaltenen Bausteine.

\textbf{Ebene 2} zoomt in einige Bausteine der Ebene 1 hinein. Sie
enthält somit die Whitebox-Beschreibungen ausgewählter Bausteine der
Ebene 1, jeweils zusammen mit Blackbox-Beschreibungen darin enthaltener
Bausteine.

\textbf{Ebene 3} zoomt in einige Bausteine der Ebene 2 hinein, usw.

\textbf{Weiterführende Informationen}

Siehe \href{https://docs.arc42.org/section-5/}{Bausteinsicht} in der
online-Dokumentation (auf Englisch!).

\subsection{Whitebox Gesamtsystem}\label{_whitebox_gesamtsystem}

An dieser Stelle beschreiben Sie die Zerlegung des Gesamtsystems anhand
des nachfolgenden Whitebox-Templates. Dieses enthält:

\begin{itemize}
\item
  Ein Übersichtsdiagramm
\item
  die Begründung dieser Zerlegung
\item
  Blackbox-Beschreibungen der hier enthaltenen Bausteine. Dafür haben
  Sie verschiedene Optionen:

  \begin{itemize}
  \item
    in \emph{einer} Tabelle, gibt einen kurzen und pragmatischen
    Überblick über die enthaltenen Bausteine sowie deren Schnittstellen.
  \item
    als Liste von Blackbox-Beschreibungen der Bausteine, gemäß dem
    Blackbox-Template (siehe unten). Diese Liste können Sie, je nach
    Werkzeug, etwa in Form von Unterkapiteln (Text), Unter-Seiten (Wiki)
    oder geschachtelten Elementen (Modellierungswerkzeug) darstellen.
  \end{itemize}
\item
  (optional:) wichtige Schnittstellen, die nicht bereits im
  Blackbox-Template eines der Bausteine erläutert werden, aber für das
  Verständnis der Whitebox von zentraler Bedeutung sind. Aufgrund der
  vielfältigen Möglichkeiten oder Ausprägungen von Schnittstellen geben
  wir hierzu kein weiteres Template vor. Im schlimmsten Fall müssen Sie
  Syntax, Semantik, Protokolle, Fehlerverhalten, Restriktionen,
  Versionen, Qualitätseigenschaften, notwendige Kompatibilitäten und
  vieles mehr spezifizieren oder beschreiben. Im besten Fall kommen Sie
  mit Beispielen oder einfachen Signaturen zurecht.
\end{itemize}

\emph{\textbf{\textless Übersichtsdiagramm\textgreater{}}}

\begin{description}
\item[Begründung]
\emph{\textless Erläuternder Text\textgreater{}}
\item[Enthaltene Bausteine]
\emph{\textless Beschreibung der enthaltenen Bausteine
(Blackboxen)\textgreater{}}
\item[Wichtige Schnittstellen]
\emph{\textless Beschreibung wichtiger Schnittstellen\textgreater{}}
\end{description}

Hier folgen jetzt Erläuterungen zu Blackboxen der Ebene 1.

Falls Sie die tabellarische Beschreibung wählen, so werden Blackboxen
darin nur mit Name und Verantwortung nach folgendem Muster beschrieben:

\begin{longtable}[]{@{}
  >{\raggedright\arraybackslash}p{(\columnwidth - 2\tabcolsep) * \real{0.3333}}
  >{\raggedright\arraybackslash}p{(\columnwidth - 2\tabcolsep) * \real{0.6667}}@{}}
\toprule\noalign{}
\begin{minipage}[b]{\linewidth}\raggedright
\textbf{Name}
\end{minipage} & \begin{minipage}[b]{\linewidth}\raggedright
\textbf{Verantwortung}
\end{minipage} \\
\midrule\noalign{}
\endhead
\bottomrule\noalign{}
\endlastfoot
\emph{\textless Blackbox 1\textgreater{}} &
~\emph{\textless Text\textgreater{}} \\
\emph{\textless Blackbox 2\textgreater{}} &
~\emph{\textless Text\textgreater{}} \\
\end{longtable}

Falls Sie die ausführliche Liste von Blackbox-Beschreibungen wählen,
beschreiben Sie jede wichtige Blackbox in einem eigenen
Blackbox-Template. Dessen Überschrift ist jeweils der Namen dieser
Blackbox.

\subsubsection{\textless Name Blackbox
1\textgreater{}}\label{_name_blackbox_1}

Beschreiben Sie die \textless Blackbox 1\textgreater{} anhand des
folgenden Blackbox-Templates:

\begin{itemize}
\item
  Zweck/Verantwortung
\item
  Schnittstelle(n), sofern diese nicht als eigenständige Beschreibungen
  herausgezogen sind. Hierzu gehören eventuell auch Qualitäts- und
  Leistungsmerkmale dieser Schnittstelle.
\item
  (Optional) Qualitäts-/Leistungsmerkmale der Blackbox, beispielsweise
  Verfügbarkeit, Laufzeitverhalten o. Ä.
\item
  (Optional) Ablageort/Datei(en)
\item
  (Optional) Erfüllte Anforderungen, falls Sie Traceability zu
  Anforderungen benötigen.
\item
  (Optional) Offene Punkte/Probleme/Risiken
\end{itemize}

\emph{\textless Zweck/Verantwortung\textgreater{}}

\emph{\textless Schnittstelle(n)\textgreater{}}

\emph{\textless(Optional) Qualitäts-/Leistungsmerkmale\textgreater{}}

\emph{\textless(Optional) Ablageort/Datei(en)\textgreater{}}

\emph{\textless(Optional) Erfüllte Anforderungen\textgreater{}}

\emph{\textless(optional) Offene Punkte/Probleme/Risiken\textgreater{}}

\subsubsection{\textless Name Blackbox
2\textgreater{}}\label{_name_blackbox_2}

\emph{\textless Blackbox-Template\textgreater{}}

\subsubsection{\textless Name Blackbox
n\textgreater{}}\label{_name_blackbox_n}

\emph{\textless Blackbox-Template\textgreater{}}

\subsubsection{\textless Name Schnittstelle
1\textgreater{}}\label{_name_schnittstelle_1}

\ldots\hspace{0pt}

\subsubsection{\textless Name Schnittstelle
m\textgreater{}}\label{_name_schnittstelle_m}

\subsection{Ebene 2}\label{_ebene_2}

Beschreiben Sie den inneren Aufbau (einiger) Bausteine aus Ebene 1 als
Whitebox.

Welche Bausteine Ihres Systems Sie hier beschreiben, müssen Sie selbst
entscheiden. Bitte stellen Sie dabei Relevanz vor Vollständigkeit.
Skizzieren Sie wichtige, überraschende, riskante, komplexe oder
besonders volatile Bausteine. Normale, einfache oder standardisierte
Teile sollten Sie weglassen.

\subsubsection{\texorpdfstring{Whitebox \emph{\textless Baustein
1\textgreater{}}}{Whitebox \textless Baustein 1\textgreater{}}}\label{_whitebox_baustein_1}

\ldots\hspace{0pt}zeigt das Innenleben von \emph{Baustein 1}.

\emph{\textless Whitebox-Template\textgreater{}}

\subsubsection{\texorpdfstring{Whitebox \emph{\textless Baustein
2\textgreater{}}}{Whitebox \textless Baustein 2\textgreater{}}}\label{_whitebox_baustein_2}

\emph{\textless Whitebox-Template\textgreater{}}

\ldots\hspace{0pt}

\subsubsection{\texorpdfstring{Whitebox \emph{\textless Baustein
m\textgreater{}}}{Whitebox \textless Baustein m\textgreater{}}}\label{_whitebox_baustein_m}

\emph{\textless Whitebox-Template\textgreater{}}

\subsection{Ebene 3}\label{_ebene_3}

Beschreiben Sie den inneren Aufbau (einiger) Bausteine aus Ebene 2 als
Whitebox.

Bei tieferen Gliederungen der Architektur kopieren Sie diesen Teil von
arc42 für die weiteren Ebenen.

\subsubsection{Whitebox \textless\_Baustein
x.1\_\textgreater{}}\label{_whitebox_baustein_x_1}

\ldots\hspace{0pt}zeigt das Innenleben von \emph{Baustein x.1}.

\emph{\textless Whitebox-Template\textgreater{}}

\subsubsection{Whitebox \textless\_Baustein
x.2\_\textgreater{}}\label{_whitebox_baustein_x_2}

\emph{\textless Whitebox-Template\textgreater{}}

\subsubsection{Whitebox \textless\_Baustein
y.1\_\textgreater{}}\label{_whitebox_baustein_y_1}

\emph{\textless Whitebox-Template\textgreater{}}
\newpage
\section{Laufzeitsicht}\label{section-runtime-view}

\textbf{Inhalt}

Diese Sicht erklärt konkrete Abläufe und Beziehungen zwischen Bausteinen
in Form von Szenarien aus den folgenden Bereichen:

\begin{itemize}
\item
  Wichtige Abläufe oder \emph{Features}: Wie führen die Bausteine der
  Architektur die wichtigsten Abläufe durch?
\item
  Interaktionen an kritischen externen Schnittstellen: Wie arbeiten
  Bausteine mit Nutzern und Nachbarsystemen zusammen?
\item
  Betrieb und Administration: Inbetriebnahme, Start, Stop.
\item
  Fehler- und Ausnahmeszenarien
\end{itemize}

Anmerkung: Das Kriterium für die Auswahl der möglichen Szenarien (d.h.
Abläufe) des Systems ist deren Architekturrelevanz. Es geht nicht darum,
möglichst viele Abläufe darzustellen, sondern eine angemessene Auswahl
zu dokumentieren.

\textbf{Motivation}

Sie sollten verstehen, wie (Instanzen von) Bausteine(n) Ihres Systems
ihre jeweiligen Aufgaben erfüllen und zur Laufzeit miteinander
kommunizieren.

Nutzen Sie diese Szenarien in der Dokumentation hauptsächlich für eine
verständlichere Kommunikation mit denjenigen Stakeholdern, die die
statischen Modelle (z.B. Bausteinsicht, Verteilungssicht) weniger
verständlich finden.

\textbf{Form}

Für die Beschreibung von Szenarien gibt es zahlreiche
Ausdrucksmöglichkeiten. Nutzen Sie beispielsweise:

\begin{itemize}
\item
  Nummerierte Schrittfolgen oder Aufzählungen in Umgangssprache
\item
  Aktivitäts- oder Flussdiagramme
\item
  Sequenzdiagramme
\item
  BPMN (Geschäftsprozessmodell und -notation) oder EPKs
  (Ereignis-Prozessketten)
\item
  Zustandsautomaten
\item
  \ldots\hspace{0pt}
\end{itemize}

\textbf{Weiterführende Informationen}

Siehe \href{https://docs.arc42.org/section-6/}{Laufzeitsicht} in der
online-Dokumentation (auf Englisch!).

\subsection{\texorpdfstring{\emph{\textless Bezeichnung Laufzeitszenario
1\textgreater{}}}{\textless Bezeichnung Laufzeitszenario 1\textgreater{}}}\label{_bezeichnung_laufzeitszenario_1}

\begin{itemize}
\item
  \textless hier Laufzeitdiagramm oder Ablaufbeschreibung
  einfügen\textgreater{}
\item
  \textless hier Besonderheiten bei dem Zusammenspiel der Bausteine in
  diesem Szenario erläutern\textgreater{}
\end{itemize}

\subsection{\texorpdfstring{\emph{\textless Bezeichnung Laufzeitszenario
2\textgreater{}}}{\textless Bezeichnung Laufzeitszenario 2\textgreater{}}}\label{_bezeichnung_laufzeitszenario_2}

\ldots\hspace{0pt}

\subsection{\texorpdfstring{\emph{\textless Bezeichnung Laufzeitszenario
n\textgreater{}}}{\textless Bezeichnung Laufzeitszenario n\textgreater{}}}\label{_bezeichnung_laufzeitszenario_n}

\ldots\hspace{0pt}
\newpage
\section{Verteilungssicht}\label{section-deployment-view}

\textbf{Inhalt}

Die Verteilungssicht beschreibt:

\begin{enumerate}
\def\labelenumi{\arabic{enumi}.}
\item
  die technische Infrastruktur, auf der Ihr System ausgeführt wird, mit
  Infrastrukturelementen wie Standorten, Umgebungen, Rechnern,
  Prozessoren, Kanälen und Netztopologien sowie sonstigen Bestandteilen,
  und
\item
  die Abbildung von (Software-)Bausteinen auf diese Infrastruktur.
\end{enumerate}

Häufig laufen Systeme in unterschiedlichen Umgebungen, beispielsweise
Entwicklung-/Test- oder Produktionsumgebungen. In solchen Fällen sollten
Sie alle relevanten Umgebungen aufzeigen.

Nutzen Sie die Verteilungssicht insbesondere dann, wenn Ihre Software
auf mehr als einem Rechner, Prozessor, Server oder Container abläuft
oder Sie Ihre Hardware sogar selbst konstruieren.

Aus Softwaresicht genügt es, auf die Aspekte zu achten, die für die
Softwareverteilung relevant sind. Insbesondere bei der
Hardwareentwicklung kann es notwendig sein, die Infrastruktur mit
beliebigen Details zu beschreiben.

\textbf{Motivation}

Software läuft nicht ohne Infrastruktur. Diese zugrundeliegende
Infrastruktur beeinflusst Ihr System und/oder querschnittliche
Lösungskonzepte, daher müssen Sie diese Infrastruktur kennen.

\textbf{Form}

Das oberste Verteilungsdiagramm könnte bereits in Ihrem technischen
Kontext enthalten sein, mit Ihrer Infrastruktur als EINE Blackbox. Jetzt
zoomen Sie in diese Infrastruktur mit weiteren Verteilungsdiagrammen
hinein:

\begin{itemize}
\item
  Die UML stellt mit Verteilungsdiagrammen (Deployment diagrams) eine
  Diagrammart zur Verfügung, um diese Sicht auszudrücken. Nutzen Sie
  diese, evtl. auch geschachtelt, wenn Ihre Verteilungsstruktur es
  verlangt.
\item
  Falls Ihre Infrastruktur-Stakeholder andere Diagrammarten bevorzugen,
  die beispielsweise Prozessoren und Kanäle zeigen, sind diese hier
  ebenfalls einsetzbar.
\end{itemize}

\textbf{Weiterführende Informationen}

Siehe \href{https://docs.arc42.org/section-7/}{Verteilungssicht} in der
online-Dokumentation (auf Englisch!).

\subsection{Infrastruktur Ebene 1}\label{_infrastruktur_ebene_1}

An dieser Stelle beschreiben Sie (als Kombination von Diagrammen mit
Tabellen oder Texten):

\begin{itemize}
\item
  die Verteilung des Gesamtsystems auf mehrere Standorte, Umgebungen,
  Rechner, Prozessoren o. Ä., sowie die physischen Verbindungskanäle
  zwischen diesen,
\item
  wichtige Begründungen für diese Verteilungsstruktur,
\item
  Qualitäts- und/oder Leistungsmerkmale dieser Infrastruktur,
\item
  Zuordnung von Softwareartefakten zu Bestandteilen der Infrastruktur
\end{itemize}

Für mehrere Umgebungen oder alternative Deployments kopieren Sie diesen
Teil von arc42 für alle wichtigen Umgebungen/Varianten.

\emph{\textbf{\textless Übersichtsdiagramm\textgreater{}}}

\begin{description}
\item[Begründung]
\emph{\textless Erläuternder Text\textgreater{}}
\item[Qualitäts- und/oder Leistungsmerkmale]
\emph{\textless Erläuternder Text\textgreater{}}
\item[Zuordnung von Bausteinen zu Infrastruktur]
\emph{\textless Beschreibung der Zuordnung\textgreater{}}
\end{description}

\subsection{Infrastruktur Ebene 2}\label{_infrastruktur_ebene_2}

An dieser Stelle können Sie den inneren Aufbau (einiger)
Infrastrukturelemente aus Ebene 1 beschreiben.

Für jedes Infrastrukturelement kopieren Sie die Struktur aus Ebene 1.

\subsubsection{\texorpdfstring{\emph{\textless Infrastrukturelement
1\textgreater{}}}{\textless Infrastrukturelement 1\textgreater{}}}\label{_infrastrukturelement_1}

\emph{\textless Diagramm + Erläuterungen\textgreater{}}

\subsubsection{\texorpdfstring{\emph{\textless Infrastrukturelement
2\textgreater{}}}{\textless Infrastrukturelement 2\textgreater{}}}\label{_infrastrukturelement_2}

\emph{\textless Diagramm + Erläuterungen\textgreater{}}

\ldots\hspace{0pt}

\subsubsection{\texorpdfstring{\emph{\textless Infrastrukturelement
n\textgreater{}}}{\textless Infrastrukturelement n\textgreater{}}}\label{_infrastrukturelement_n}

\emph{\textless Diagramm + Erläuterungen\textgreater{}}
\newpage
\section{Querschnittliche Konzepte}\label{section-concepts}

\textbf{Inhalt}

Dieser Abschnitt beschreibt übergreifende, prinzipielle Regelungen und
Lösungsansätze, die an mehreren Stellen (=\emph{querschnittlich})
relevant sind.

Solche Konzepte betreffen oft mehrere Bausteine. Dazu können vielerlei
Themen gehören, wie beispielsweise die Themen aus dem nachfolgenden
Diagramm:

\begin{figure}
\centering
\includegraphics{images/08-concepts-DE.drawio.png}
\caption{}
\end{figure}

\textbf{Motivation}

Konzepte bilden die Grundlage für \emph{konzeptionelle Integrität}
(Konsistenz, Homogenität) der Architektur und damit eine wesentliche
Grundlage für die innere Qualität Ihrer Systeme.

Dieser Abschnitt im Template ist der richtige Ort für die konsistente
Behandlung solcher Themen.

Viele solche Konzepte beeinflussen oder beziehen sich auf mehrerer Ihrer
Bausteine.

\textbf{Form}

Kann vielfältig sein:

\begin{itemize}
\item
  Konzeptpapiere mit beliebiger Gliederung,
\item
  beispielhafte Implementierung speziell für technische Konzepte,
\item
  übergreifende Modelle/Szenarien mit Notationen, die Sie auch in den
  Architektursichten nutzen,
\end{itemize}

\textbf{Struktur}

Wählen Sie \textbf{nur} die wichtigsten Themen für Ihr System und
erklären das jeweilige Konzept dann unter einer Level-2 Überschrift
dieser Sektion (z.B. 8.1, 8.2 etc).

Beschränken Sie sich auf die wichtigen, und versuchen \textbf{auf keinen
Fall} alle oben dargestellten Themen zu bearbeiten.

\textbf{Weiterführende Informationen}

Einige Themen innerhalb von Systemen betreffen oft mehrere Bausteine,
Hardwareelemente oder Prozesse. Es könnte einfacher sein, solche
\emph{Querschnittsthemen} an einer zentralen Stelle zu kommunizieren
oder zu dokumentieren, anstatt sie in der Beschreibung der betreffenden
Bausteine, Hardwareelemente oder Entwicklungsprozesse zu wiederholen.

Bestimmte Konzepte können \textbf{alle} Elemente eines Systems
betreffen, andere sind vielleicht nur für einige wenige relevant.

Siehe \href{https://docs.arc42.org/section-8/}{Querschnittliche
Konzepte} in der online-Dokumentation (auf Englisch).

\subsection{\texorpdfstring{\emph{\textless Konzept
1\textgreater{}}}{\textless Konzept 1\textgreater{}}}\label{_konzept_1}

\emph{\textless Erklärung\textgreater{}}

\subsection{\texorpdfstring{\emph{\textless Konzept
2\textgreater{}}}{\textless Konzept 2\textgreater{}}}\label{_konzept_2}

\emph{\textless Erklärung\textgreater{}}

\ldots\hspace{0pt}

\subsection{\texorpdfstring{\emph{\textless Konzept
n\textgreater{}}}{\textless Konzept n\textgreater{}}}\label{_konzept_n}

\emph{\textless Erklärung\textgreater{}}
\newpage
\section{Architekturentscheidungen}\label{section-design-decisions}

\textbf{Inhalt}

Wichtige, teure, große oder riskante Architektur- oder
Entwurfsentscheidungen inklusive der jeweiligen Begründungen. Mit
"Entscheidungen" meinen wir hier die Auswahl einer von mehreren
Alternativen unter vorgegebenen Kriterien.

Wägen Sie ab, inwiefern Sie Entscheidungen hier zentral beschreiben,
oder wo eine lokale Beschreibung (z.B. in der Whitebox-Sicht von
Bausteinen) sinnvoller ist. Vermeiden Sie Redundanz. Verweisen Sie evtl.
auf Abschnitt 4, wo schon grundlegende strategische Entscheidungen
beschrieben wurden.

\textbf{Motivation}

Stakeholder des Systems sollten wichtige Entscheidungen verstehen und
nachvollziehen können.

\textbf{Form}

Verschiedene Möglichkeiten:

\begin{itemize}
\item
  ADR
  (\href{https://cognitect.com/blog/2011/11/15/documenting-architecture-decisions}{Documenting
  Architecture Decisions}) für jede wichtige Entscheidung
\item
  Liste oder Tabelle, nach Wichtigkeit und Tragweite der Entscheidungen
  geordnet
\item
  ausführlicher in Form einzelner Unterkapitel je Entscheidung
\end{itemize}

\textbf{Weiterführende Informationen}

Siehe
\href{https://docs.arc42.org/section-9/}{Architekturentscheidungen} in
der arc42 Dokumentation (auf Englisch!). Dort finden Sie Links und
Beispiele zum Thema ADR.
\newpage
\section{Qualitätsanforderungen}\label{section-quality-scenarios}

\textbf{Inhalt}

Dieser Abschnitt enthält alle relevanten Qualitätsanforderungen.

Die wichtigsten davon haben Sie bereits in Abschnitt 1.2
(Qualitätsziele) hervorgehoben, daher soll hier nur auf sie verwiesen
werden. In diesem Abschnitt 10 sollten Sie auch Qualitätsanforderungen
mit geringerer Bedeutung erfassen, deren Nichterfüllung keine großen
Risiken birgt (die aber \emph{nice-to-have} sein könnten).

\textbf{Motivation}

Weil Qualitätsanforderungen die Architekturentscheidungen oft maßgeblich
beeinflussen, sollten Sie die für Ihre Stakeholder relevanten
Qualitätsanforderungen kennen, möglichst konkret und operationalisiert.

\begin{itemize}
\item
  Siehe
  \href{https://docs.arc42.org/section-10/}{Qualitätsanforderungen} in
  der online-Dokumentation (auf Englisch!).
\item
  Siehe auch das ausführliche \href{https://quality.arc42.org}{Q42
  Qualitätsmodell auf https://quality.arc42.org}.
\end{itemize}

\subsection{Übersicht der
Qualitätsanforderungen}\label{_uxfcbersicht_der_qualituxe4tsanforderungen}

\textbf{Inhalt}

Eine Übersicht oder Zusammenfassung der Qualitätsanforderungen.

\textbf{Motivation}

Oft stößt man auf Dutzende (oder sogar Hunderte) von detaillierten
Qualitätsanforderungen für ein System. In diesem Abschnitt sollten Sie
versuchen, sie zusammenzufassen, z. B. durch die Beschreibung von
Kategorien oder Themen (wie z.B. von
\href{https://www.iso.org/obp/ui/\#iso:std:iso-iec:25010:ed-2:v1:en}{ISO
25010:2023} oder \href{https://quality.arc42.org}{Q42} vorgeschlagen).

Wenn diese Kurzbeschreibungen oder Zusammenfassungen bereits präzise,
spezifisch und messbar sind, können Sie Abschnitt 10.2 auslassen.

\textbf{Form}

Verwenden Sie eine einfache Tabelle, in der jede Zeile eine Kategorie
oder ein Thema und eine kurze Beschreibung der Qualitätsanforderung
enthält. Alternativ können Sie auch eine Mindmap verwenden, um diese
Qualitätsanforderungen zu strukturieren. In der Literatur (insb.
{[}Bass+21{]}) ist die Idee eines \emph{Quality Attribute Utility Tree}
(auf Deutsch manchmal kurz als \emph{Qualitätsbaum} bezeichnet)
beschrieben worden, der den Oberbegriff „Qualität`` als Wurzel hat und
eine baumartige Verfeinerung des Begriffs „Qualität`` verwendet.

\subsection{Qualitätsszenarien}\label{_qualituxe4tsszenarien}

\textbf{Inhalt}

Qualitätsszenarien konkretisieren Qualitätsanforderungen und ermöglichen
es zu entscheiden, ob sie erfüllt sind (im Sinne von
Akzeptanzkriterien). Stellen Sie sicher, dass Ihre Szenarien spezifisch
und messbar sind.

Zwei Arten von Szenarien finden wir besonders nützlich:

\begin{itemize}
\item
  Nutzungsszenarien (auch bekannt als Anwendungs- oder
  Anwendungsfallszenarien) beschreiben, wie das System zur Laufzeit auf
  einen bestimmten Auslöser reagieren soll. Hierunter fallen auch
  Szenarien zur Beschreibung von Effizienz oder Performance. Beispiel:
  Das System beantwortet eine Benutzeranfrage innerhalb einer Sekunde.
\item
  Änderungsszenarien\_ beschreiben die gewünschte Wirkung einer Änderung
  oder Erweiterung des Systems oder seiner unmittelbaren Umgebung.
  Beispiel: Zusätzliche Funktionalität wird implementiert oder
  Anforderungen an ein Qualitätsmerkmal ändern sich, und der Aufwand
  oder die Dauer der Änderung wird gemessen.
\end{itemize}

\textbf{Form}

Typische Informationen für detaillierte Szenarien sind die folgenden:

In Kurzform (bevorzugt im Q42-Modell):

\begin{itemize}
\item
  K\textbf{ontext/Hintergrund}: Um welche Art von System oder Komponente
  handelt es sich, wie sieht die Umgebung oder Situation aus?
\item
  \textbf{Quelle/Stimulus}: Wer oder was initiiert oder löst ein
  Verhalten, eine Reaktion oder eine Aktion aus.
\item
  \textbf{Metrik/Akzeptanzkriterien}: Eine Reaktion einschließlich einer
  \emph{Maßnahme} oder \emph{Metrik}
\end{itemize}

Die Langform von Szenarien (die von der SEI und {[}Bass+21{]} bevorzugt
wird) ist detaillierter und enthält die folgenden Informationen:

\begin{itemize}
\item
  \textbf{Szenario-ID}: Ein eindeutiger Bezeichner für das Szenario.
\item
  \textbf{Szenario-Name}: Ein kurzer, beschreibender Name für das
  Szenario.
\item
  \textbf{Quelle}: Die Entität (Benutzer, System oder Ereignis), die das
  Szenario auslöst.
\item
  \textbf{Stimulus}: Das auslösende Ereignis oder die Bedingung, auf die
  das System reagieren muss.
\item
  \textbf{Umgebung}: Der betriebliche Kontext oder die Bedingungen,
  unter denen das System den Stimulus erlebt.
\item
  \textbf{Artefakt}: Die Bausteine oder anderen Elemente des Systems,
  die von dem Stimulus betroffen sind.
\item
  \textbf{Reaktion}: Das Ergebnis oder Verhalten, das das System als
  Reaktion auf den Stimulus zeigt.
\item
  \textbf{Antwortmaß}: Das Kriterium oder die Metrik, nach der die
  Antwort des Systems bewertet wird.
\end{itemize}

\textbf{Beispiele}

Ausführliche Beispiele für Qualitätsanforderungen finden Sie auf
\href{https://quality.arc42.org}{der Website zum Qualitätsmodell Q42}.

\begin{itemize}
\item
  Len Bass, Paul Clements, Rick Kazman: „Software Architecture in
  Practice``, 4. Auflage, Addison-Wesley, 2021.
\end{itemize}
\newpage
\section{Risiken und technische Schulden}\label{section-technical-risks}

\textbf{Inhalt}

Eine nach Prioritäten geordnete Liste der erkannten Architekturrisiken
und/oder technischen Schulden.

\begin{quote}
Risikomanagement ist Projektmanagement für Erwachsene.

---  Tim Lister Atlantic Systems Guild
\end{quote}

Unter diesem Motto sollten Sie Architekturrisiken und/oder technische
Schulden gezielt ermitteln, bewerten und Ihren Management-Stakeholdern
(z.B. Projektleitung, Product-Owner) transparent machen.

\textbf{Form}

Liste oder Tabelle von Risiken und/oder technischen Schulden, eventuell
mit vorgeschlagenen Maßnahmen zur Risikovermeidung, Risikominimierung
oder dem Abbau der technischen Schulden.

\textbf{Weiterführende Informationen}

Siehe \href{https://docs.arc42.org/section-11/}{Risiken und technische
Schulden} in der online-Dokumentation (auf Englisch!).

\section{Glossar}\label{section-glossary}

\textbf{Inhalt}

Die wesentlichen fachlichen und technischen Begriffe, die Stakeholder im
Zusammenhang mit dem System verwenden.

Nutzen Sie das Glossar ebenfalls als Übersetzungsreferenz, falls Sie in
mehrsprachigen Teams arbeiten.

\textbf{Motivation}

Sie sollten relevante Begriffe klar definieren, so dass alle Beteiligten

\begin{itemize}
\item
  diese Begriffe identisch verstehen, und
\item
  vermeiden, mehrere Begriffe für die gleiche Sache zu haben.
\end{itemize}

\textbf{Form}

Zweispaltige Tabelle mit \textless Begriff\textgreater{} und
\textless Definition\textgreater.

Eventuell weitere Spalten mit Übersetzungen, falls notwendig.

\textbf{Weiterführende Informationen}

Siehe \href{https://docs.arc42.org/section-12/}{Glossar} in der
online-Dokumentation (auf Englisch!).

\begin{longtable}[]{@{}
  >{\raggedright\arraybackslash}p{(\columnwidth - 2\tabcolsep) * \real{0.3333}}
  >{\raggedright\arraybackslash}p{(\columnwidth - 2\tabcolsep) * \real{0.6667}}@{}}
\toprule\noalign{}
\begin{minipage}[b]{\linewidth}\raggedright
Begriff
\end{minipage} & \begin{minipage}[b]{\linewidth}\raggedright
Definition
\end{minipage} \\
\midrule\noalign{}
\endhead
\bottomrule\noalign{}
\endlastfoot
\emph{\textless Begriff-1\textgreater{}} &
\emph{\textless Definition-1\textgreater{}} \\
\emph{\textless Begriff-2} &
\emph{\textless Definition-2\textgreater{}} \\
\end{longtable}

\end{document}